% ===========================================================================
% DriftDream -- Introduction, Related Work, Method
% Drop-in sections. Replace the preamble with the track template when
% submitting; nothing below \begin{document} depends on the class.
% ===========================================================================
\documentclass[10pt]{article}

\usepackage[margin=1in]{geometry}
\usepackage{amsmath,amssymb,amsthm}
\usepackage{booktabs}
\usepackage{graphicx}
\usepackage{algorithm}
\usepackage{algpseudocode}
\usepackage[colorlinks=true,allcolors=blue]{hyperref}
\usepackage[T1]{fontenc}
\usepackage[expansion=false]{microtype}
\usepackage[round,authoryear]{natbib}

% --- notation --------------------------------------------------------------
\newcommand{\R}{\mathbb{R}}
\newcommand{\bx}{\mathbf{x}}
\newcommand{\bp}{\mathbf{p}}
\newcommand{\bc}{\mathbf{c}}
\newcommand{\bd}{\boldsymbol{\delta}}
\newcommand{\bn}{\mathbf{n}}
\newcommand{\bu}{\mathbf{u}}
\newcommand{\bv}{\mathbf{v}}
\newcommand{\bF}{\mathbf{F}}
\newcommand{\bff}{\mathbf{f}}
\newcommand{\bA}{\mathbf{A}}
\newcommand{\bb}{\mathbf{b}}
\newcommand{\W}{\mathcal{W}}
\newcommand{\Loss}{\mathcal{L}}
\newcommand{\Om}{\Omega}

\title{DriftDream: Training-Free Generation of Static Images\\
       with a Prescribed Illusory-Motion Field}
\author{}
\date{}

\begin{document}
\maketitle

% ===========================================================================
\section{Introduction}
% ===========================================================================

A still image can look like it is moving. Patterns in the peripheral drift
family, of which Kitaoka's \emph{Rotating Snakes} is the best known example,
are printed on paper and do not change, yet most viewers see the rings turn
\cite{faubert1999peripheral,kitaoka2003snakes}. The effect comes from repeating
elements whose luminances follow an asymmetric order, together with the small
involuntary eye movements that occur during fixation. Elements at different
luminances are registered by the visual system at slightly different times, and
a motion detector reads that timing difference as movement
\cite{conway2005neural,backus2005illusory}. The percept is stronger in
observers whose fixation is less stable, which ties it directly to eye
movements rather than to the image alone \cite{murakami2006fixation}.

Recent work has shown that a frozen text-to-image diffusion model can be made
to produce illusions without any training, by splitting the denoising signal
along an axis on which one canvas can satisfy two descriptions at once. Visual
Anagrams uses pixel permutations, so the image shows one thing upright and
another when flipped or rotated \cite{geng2024anagrams}. Factorized Diffusion
splits by spatial frequency, giving hybrid images in the sense of
\citet{oliva2006hybrid} \cite{geng2024factorized}. Images that Sound splits
across modality, so a single canvas is both a picture and an audio spectrogram
\cite{chen2024imagesthatsound}.

Every one of these axes is instantaneous. The image is read once, and the two
readings are of the same frozen array of pixels. Motion is not on the list, and
the reason is not that nobody thought of it. It is that the obvious way to
write the objective does not work.

\paragraph{Why the obvious objective is empty.}
Suppose we want a still image whose perceived motion matches a target field
$\bF^\ast$, and we have a differentiable optical-flow model $M$. The natural
loss is to read the flow of the image and compare it to $\bF^\ast$. But a still
image has no flow. Feed $M$ the same frame twice and it returns approximately
zero everywhere, for every image, so the loss is constant and its gradient is
zero. There is nothing to descend. The illusion we are trying to create does
not live in a single frame; it appears only when the image moves slightly on
the retina.

\paragraph{What we do instead.}
We put the eye movements into the objective. At each optimisation step we take
the current image, apply a small set of rigid shifts that stand in for
microsaccades, read the flow between consecutive shifted frames, remove the
part of that flow explained by the shift itself, and average what is left over
several shift directions. The average residual is the illusory component. We
compare it to $\bF^\ast$ with a normalised correlation and use the result as a
guidance term on a frozen generator. We call this \emph{jitter-ensemble motion
guidance}. It is the one new piece of machinery in the paper, and it is what
turns an empty objective into a signal that can be followed.

Two properties follow from the construction rather than from tuning. The
residual is signed, so asking for clockwise and asking for anticlockwise give
different images rather than the same image twice. And it is defined per pixel,
so the target can be localised: rotation inside one ring, or drift confined to
the shape of a letter while the surround stays still.

\paragraph{What we do not do.}
No network is trained, fine-tuned, or adapted. Every model in the pipeline is
frozen and used only in the forward and backward pass. The optimised variables
are the image (or the diffusion latent), nothing else.

\paragraph{Why the evaluation is the hard part.}
A model that scores an image is also a model that can be fooled. Optimising
pixels against a flow network is close to the setting in which adversarial
examples are found, and an image that drives $M$ to any chosen output need not
look like it moves to a person. This concern is not hypothetical for this
particular percept. \citet{rosario2026illusory} tested a range of flow models
on Rotating Snakes and found that most produce flow fields that do not match
what people report; only a model built to mirror the V1--MT pathway
\cite{sun2025multiorder} showed the expected rotation, and only when the
stimulus was jittered to imitate eye movements. On the predictive-coding side,
\citet{watanabe2018prednet} reported that a PredNet trained only to predict the
next frame reproduces the illusion, but \citet{kirubeswaran2023inconsistent}
retrained the same architecture ten times and found the behaviour came and went
across runs, and never appeared for greyscale versions that people still see.

We therefore separate the model that is optimised from the models that judge.
Guidance uses one reader. Reporting uses readers that never appear in any
gradient, plus a two-alternative forced-choice study in which people report
only the direction they see. We also require a negative control: a transform
that destroys the percept for human viewers must drive the measured correlation
to zero. Building that control turned out to be less obvious than expected, and
we describe in Section~\ref{sec:controls} why the first version we tried made
the score go up rather than down.

\paragraph{Contributions.}
\begin{enumerate}\itemsep2pt
\item A training-free method that generates a still image whose \emph{perceived}
      motion follows a field the user specifies, including its direction and
      where in the frame it appears. Prior illusion work with diffusion priors
      is instantaneous; prior work on generating motion illusions
      \cite{sinapayen2021eigen} controls neither content nor direction.
\item Jitter-ensemble motion guidance: an objective defined on the residual
      flow averaged over simulated microsaccades in several directions. We show
      the ensemble is required, not cosmetic --- with zero jitter the measured
      correlation is zero by construction, and with one direction it is
      dominated by reader artefacts.
\item An evaluation protocol built around held-out readers, a signed control
      (asking for the opposite direction must reverse the measurement), and a
      negative control that must collapse the score. We report a case in which
      an intuitive control fails, and what replaces it.
\end{enumerate}

% ===========================================================================
\section{Related Work}
% ===========================================================================

\subsection{Illusions from frozen diffusion models}

\citet{geng2024anagrams} generate multi-view illusions by estimating noise from
several transformed views of the same noisy image and averaging the estimates
before each reverse step. Their analysis shows the argument holds when the
views are orthogonal transformations, which covers flips, rotations by
multiples of $90^\circ$, colour inversion, and jigsaw permutations.
\citet{geng2024factorized} replace the transformation with a decomposition of
the noise into frequency bands, recovering hybrid images
\cite{oliva2006hybrid} and related effects. \citet{chen2024imagesthatsound}
denoise one canvas under an image model and an audio-spectrogram model at the
same time, so the result can be viewed and played.

These share a structure: a fixed decomposition of the canvas, and two or more
conditions imposed on the parts. Our objective is not a decomposition. There is
no second prompt and no second view. There is one image, and a measurement made
on it by a model that is shown the image in motion. The quantity being
controlled is not what the image depicts under a transformation but what
happens in an observer when the image is held still.

\subsection{Guidance and search at inference time}

Classifier guidance \cite{dhariwal2021guidance} and classifier-free guidance
\cite{ho2022cfg} established that the sampling trajectory can be pushed toward
an objective without retraining. Universal guidance \cite{bansal2023universal}
and self-guidance \cite{epstein2023selfguidance} extend this to arbitrary
off-the-shelf networks and to the model's own internal activations. We follow
the same recipe: differentiate a loss computed on the current estimate of the
clean image and add the gradient to the state.

Closest in mechanism is Motion Guidance \cite{geng2024motionguidance}, which
backpropagates through a frozen flow network \cite{teed2020raft} to move
content in an edit. The distinction matters and is easy to blur. Their target
displacement is measured between two \emph{different} images, a source and an
output, and the content really does move. Ours is measured between an image and
a shifted copy of \emph{itself}, and nothing moves; what changes is what a
motion model reports. Applying their objective to a single still frame reduces
exactly to our zero-jitter ablation, which is why we report that ablation
instead of implementing a separate baseline.

When gradients are unavailable or badly behaved, one can search instead.
\citet{ma2025inferencescaling} search over initial noises using a verifier, and
\citet{singhal2025fk} use interacting-particle Feynman--Kac methods to steer
sampling with an arbitrary reward. Both apply to our objective without change,
since it is a scalar score, and we keep them as a fallback path.

\subsection{Peripheral drift and its explanations}

\citet{faubert1999peripheral} described the peripheral drift illusion, and
\citet{kitaoka2003snakes} characterised the Rotating Snakes variant, in which
motion is perceived following the repeating order black, dark, white, light.
\citet{conway2005neural} showed the effect survives when colour is removed, so
the critical variable is the luminance relation of neighbouring elements, and
recorded direction-selective cells in macaque V1 and MT that respond
directionally to static element pairs. They report that high-contrast elements
(black and white, relative to the mean) reach peak response 10--20\,ms earlier
than the intermediate greys, which supplies the timing difference a motion
detector needs. \citet{backus2005illusory} give an account in terms of the
time course of responses to contrast and luminance, and
\citet{fermuller2010causal} attribute the effect to causal temporal filtering.
\citet{murakami2006fixation} found that observers with less stable fixation
report stronger illusory motion, which is the empirical basis for building eye
movements into our objective rather than treating them as noise.

\subsection{Models that do or do not reproduce the illusion}

\citet{watanabe2018prednet} showed that a PredNet \cite{lotter2017prednet}
trained for next-frame prediction produces rotational flow on Rotating Snakes.
\citet{kirubeswaran2023inconsistent} retrained the architecture from ten
different initialisations and found the behaviour was not reliable across runs,
and that no instance produced rotation for greyscale variants that people do
see. \citet{sun2025multiorder} build a model of the V1--MT pathway with a
trainable motion-energy stage and a recurrent attention stage, extending an
earlier version \cite{sun2023motionenergy}. \citet{rosario2026illusory} compare
several flow models on Rotating Snakes and find that most fail; the
dual-channel model of \citet{sun2025multiorder} shows the expected rotation,
most clearly under simulated saccades. They also note that all models tested
lack any mechanism for peripheral vision, while the human effect is stronger
away from the fovea.

We take two things from this line of work and one warning. We take the
simulated-saccade protocol and the normalised flow correlation as our
measurement, and the target field of \citet{rosario2026illusory} as our
default $\bF^\ast$ (Section~\ref{sec:target}). The warning is that no single
reader can carry the claim, which is why our reporting is cross-model.

\subsection{Generating motion illusions}

\citet{sinapayen2021eigen} search for motion illusions with an evolutionary
algorithm, scoring candidates with a PredNet, and confirm with a survey that
the results work on people. That is the closest existing system, and the
differences are in what can be asked for. Their search has no content term, so
outputs are abstract repeating patterns rather than pictures of anything; it
scores an image with a scalar, so there is no way to ask for a particular
direction or to place the effect in a chosen region; and it is gradient-free
evolution over a pattern-generating network rather than guidance of a
generative prior. Our objective is a per-pixel signed field, which is what
makes direction and layout controllable, and it attaches to a diffusion model,
which is what lets the image also be a picture.

Relatedly, \citet{gomezvilla2019deceived} showed that convolutional networks
can be fooled by classical illusions, and \citet{poole2023dreamfusion}
popularised the pattern of optimising a signal through a frozen model, which we
use in a different setting.

% ===========================================================================
\section{Method}
% ===========================================================================

\subsection{Overview and notation}

The input is a target motion field $\bF^\ast:\Om\to\R^2$ defined on a support
mask $\Om\subseteq\{1..H\}\times\{1..W\}$, an optional text prompt, and a frozen
generator. The output is an image $\bx\in[0,1]^{3\times H\times W}$.

Write $\bp=(p_x,p_y)$ for a pixel with $p_y$ increasing downwards, matching
image row order. A flow field $\bff\in\R^{2\times H\times W}$ has $\bff_1$ the
horizontal and $\bff_2$ the vertical component, again with positive meaning
downwards. A reader $M$ maps a stack of frames to one flow field. We use
$m(\bp)\in[0,1]$ for the soft mask.

The pipeline has four parts, described in turn: the target field
(\S\ref{sec:target}), the simulated fixation (\S\ref{sec:fixation}), the
residual flow that forms the signal (\S\ref{sec:residual}), and the loss and
guidance step (\S\ref{sec:loss}). Section~\ref{sec:reader} describes the reader
we supply so that the method runs without any gated download, and
Section~\ref{sec:controls} the controls.

\subsection{Target field}
\label{sec:target}

For rotation we use the field of \citet{rosario2026illusory}, who introduced it
to score models against the human percept of Rotating Snakes. With centre $\bc$,
outer radius $R$, radial exponent $\gamma$, and $r=\lVert\bp-\bc\rVert$, the
unit tangent and the field are
\begin{equation}
\bu(\bp)=\frac{1}{r}\begin{pmatrix}-(p_y-c_y)\\ p_x-c_x\end{pmatrix},
\qquad
\tilde\bF(\bp)=\omega\left(\frac{r}{R}\right)^{\!\gamma}\bu(\bp)\, m(\bp).
\label{eq:rotation}
\end{equation}
The sign $\omega=\pm1$ selects the direction and $\gamma=1$ gives the linear
radial profile of rigid rotation, avoiding a magnitude discontinuity at the
outer edge. A check on the convention: at the top of the circle
$p_y<c_y$ and $p_x=c_x$, so $\bu=(1,0)$; with $p_y$ downwards, motion to the
right at the top of a circle is clockwise, so $\omega=+1$ means clockwise on
screen.

Radial expansion replaces $\bu$ with the outward unit radial vector, and
laminar drift replaces it with a constant unit vector inside a stencil,
allowing targets such as ``this word drifts left while the background is
still''.

The field is normalised over the mask,
\begin{equation}
\bF^\ast=\tilde\bF\,\big/\,\lVert\tilde\bF\rVert_{\Om},
\qquad
\lVert\bff\rVert_{\Om}^2=\sum_{\bp}m(\bp)^2\lVert\bff(\bp)\rVert^2 ,
\label{eq:normalise}
\end{equation}
because the score in \eqref{eq:rho} is scale invariant: only direction and
relative magnitude are being asked for, not a speed in pixels per second, which
a still image does not have.

\subsection{Simulated fixation}
\label{sec:fixation}

Let $\W_{\bd}$ be a translation by $\bd\in\R^2$ implemented by bilinear
sampling,
\begin{equation}
(\W_{\bd}\bx)(\bp)=\bx(\bp-\bd),
\label{eq:warp}
\end{equation}
so image content moves by $+\bd$. Sampling outside the frame uses reflection
padding; zero padding puts a hard edge at the border and the reader reports
large spurious flow there. $\W_{\bd}$ is differentiable in $\bx$, which is what
lets gradients pass from the reader back to the image.

A fixation clip is a sequence of cumulative offsets
$\bd_1=\mathbf 0,\bd_2,\dots,\bd_K$ along a straight path with step length
$\Delta$. We use $J$ such paths with directions spread evenly over the circle,
$\theta_j=2\pi j/J$. From \eqref{eq:warp}, the true flow from frame $i-1$ to
frame $i$ of path $j$ is the constant field $\bd_{j,i}-\bd_{j,i-1}$.

Using several directions is not a detail. With one direction the measurement is
dominated by whatever that particular shift excites in the reader: on a fixed
image we measured the score swinging between $+0.13$ and $-0.09$ across jitter
seeds. Averaging over directions removes most of that, and it also removes an
easy exploit, since structure tuned to one shift direction does not survive the
average.

\subsection{Residual flow}
\label{sec:residual}

The flow the reader returns is dominated by the shift we applied. The illusory
part is what remains. Two ways to remove the shift are available and they do
not behave the same.

The direct choice subtracts the known offset,
$\bff^{\text{res}}=M(\cdot)-(\bd_{j,i}-\bd_{j,i-1})$. This is only correct if
the reader's gain is one. Real motion models under-report speed. On a smoothed
random texture at $96\times96$, our built-in reader recovers a 2-pixel shift as
$0.73$ pixels, a gain of $0.37$, so the subtraction leaves a residual
proportional to the shift, pointing along the jitter, and far larger than the
effect we are looking for: with $K=4$, $\Delta=5$ and a single direction, mean
residual magnitude $4.95$ pixels.

We instead remove the uniform component the reader actually reported, using the
spatial median of each component,
\begin{equation}
\bff^{\text{res}}_{j,i}
= M\!\left(\W_{\bd_{j,i-1}}\bx,\; \W_{\bd_{j,i}}\bx\right)
- \operatorname{med}_{\bp}\,
  M\!\left(\W_{\bd_{j,i-1}}\bx,\; \W_{\bd_{j,i}}\bx\right).
\label{eq:residual}
\end{equation}
On the same test this leaves $0.33$ pixels, a factor of 15 smaller. The median
is used rather than the mean so that a strong local response does not shift the
estimate of the global component.

There is also a perceptual reason to prefer \eqref{eq:residual}. A uniform
translation of the whole retinal image is what an eye movement produces, and
observers do not perceive it as object motion. What is perceived is the part of
the field that departs from uniform. The cost is that a genuinely uniform
illusory drift across the entire frame would be invisible to the measurement.
Every target we use is spatially structured, so this does not bind, but a
whole-frame constant target would need the direct subtraction instead.

The ensemble average over frames and directions is
\begin{equation}
\bar\bff(\bx)=\frac{1}{J(K-1)}\sum_{j=1}^{J}\sum_{i=2}^{K}
\bff^{\text{res}}_{j,i}.
\label{eq:ensemble}
\end{equation}
When $K=1$ or $\Delta=0$ this is zero for every image, which is the empty
objective from the introduction, and is the ablation reported in the
experiments.

One consequence of \eqref{eq:ensemble} is worth stating because it affects how
the ablations should be read. When the $J$ directions are spread evenly, the
offsets sum to approximately zero, so the gain error described above also
cancels in the average. The direct subtraction and \eqref{eq:residual}
therefore look similar once $J$ is large, and they differ sharply only at
$J=1$. Diagnostics comparing the two must use a single direction or they
measure nothing.

\subsection{Score, loss, and guidance}
\label{sec:loss}

The score is the masked normalised correlation used by
\citet{rosario2026illusory},
\begin{equation}
\rho(\bar\bff,\bF^\ast)=
\frac{\sum_{\bp}m(\bp)^2\,\langle\bar\bff(\bp),\bF^\ast(\bp)\rangle}
     {\lVert\bar\bff\rVert_{\Om}\;\lVert\bF^\ast\rVert_{\Om}+\varepsilon}
\;\in[-1,1].
\label{eq:rho}
\end{equation}
By Cauchy--Schwarz $\rho=1$ exactly when $\bar\bff\propto\bF^\ast$ with a
positive constant on $\Om$; $\rho(\bF,-\bF)=-1$; and $\rho$ vanishes for fields
orthogonal at every pixel, such as a rotation field against its own radial
counterpart. We verified all three numerically to within $10^{-3}$.

The full objective adds a content term and, for stencil targets, a quiet term,
\begin{equation}
\Loss(\bx)=
\underbrace{\lambda_m\bigl(1-\rho(\bar\bff(\bx),\bF^\ast)\bigr)}_{\text{motion}}
+\underbrace{\lambda_c\bigl(1-\cos(E_I(\bx),E_T(\text{prompt}))\bigr)}_{\text{content}}
+\underbrace{\lambda_q\,Q(\bar\bff(\bx))}_{\text{quiet}},
\label{eq:loss}
\end{equation}
where $E_I,E_T$ are frozen CLIP image and text encoders
\cite{radford2021clip} and
\begin{equation}
Q(\bff)=\frac{\sum_{\bp}\bigl(1-m(\bp)\bigr)^{2}\lVert\bff(\bp)\rVert^{2}}
              {2\sum_{\bp}\bigl(1-m(\bp)\bigr)+\varepsilon}
\label{eq:quiet}
\end{equation}
penalises residual flow outside the support. Without $Q$ a stencil target tends
to make the whole frame shimmer, and ``only the word moves'' is not legible.

\paragraph{Guidance.} For a diffusion generator with noise prediction
$\epsilon_\theta$ and cumulative noise schedule $\bar\alpha_t$, the current
estimate of the clean image is the standard Tweedie form
\begin{equation}
\hat\bx_0(\bx_t,t)=
\frac{\bx_t-\sqrt{1-\bar\alpha_t}\;\epsilon_\theta(\bx_t,t,c)}
     {\sqrt{\bar\alpha_t}},
\label{eq:tweedie}
\end{equation}
which follows from the forward process
$\bx_t=\sqrt{\bar\alpha_t}\,\bx_0+\sqrt{1-\bar\alpha_t}\,\epsilon$
\cite{ho2020ddpm}. We evaluate \eqref{eq:loss} at $\hat\bx_0$ and take a step
against its gradient in the state, following
\cite{dhariwal2021guidance,bansal2023universal},
\begin{equation}
\bx_t\;\leftarrow\;\bx_t-s_t\,\Pi_{\tau}\!\left[
\nabla_{\bx_t}\Loss\bigl(\hat\bx_0(\bx_t,t)\bigr)\right],
\label{eq:guidance}
\end{equation}
with $\Pi_\tau$ a norm clip at $\tau$ and $s_t$ a decaying scale. Gradients
pass through the reader and the warps only; $\epsilon_\theta$ is treated as
constant within the step, so no weights are touched and no graph is kept through
the sampler. For a latent model the decoder sits between $\hat\bx_0$ and the
loss and is likewise frozen and differentiated through
\cite{rombach2022ldm,podell2024sdxl}. Guidance is applied on an interior window
of steps: early states are too noisy for a motion reader to say anything useful,
and late steps mainly add texture.

\begin{algorithm}[t]
\caption{Jitter-ensemble motion guidance (one sample)}
\label{alg:main}
\begin{algorithmic}[1]
\Require target $\bF^\ast$, mask $m$, prompt $c$, frozen generator and reader
         $M$, offsets $\{\bd_{j,i}\}$, weights $\lambda_m,\lambda_c,\lambda_q$
\State $\bx_T\sim\mathcal N(\mathbf 0,\mathbf I)$
\For{$t=T,\dots,1$}
  \State $\hat\bx_0\gets$ Tweedie estimate \eqref{eq:tweedie}
  \If{$t$ in guidance window}
    \State $\bar\bff\gets$ ensemble residual \eqref{eq:residual}--\eqref{eq:ensemble}
           evaluated at $\hat\bx_0$
    \State $\Loss\gets\lambda_m(1-\rho(\bar\bff,\bF^\ast))
            +\lambda_c\Loss_{\text{content}}+\lambda_q Q(\bar\bff)$
    \State $\bx_t\gets\bx_t-s_t\,\Pi_\tau[\nabla_{\bx_t}\Loss]$
  \EndIf
  \State $\bx_{t-1}\gets\textsc{SamplerStep}(\bx_t,\epsilon_\theta,t)$
\EndFor
\State \Return $\bx_0$
\end{algorithmic}
\end{algorithm}

Because \eqref{eq:loss} is a scalar function of the image, the same objective
drives gradient-free optimisers when gradients are unavailable: best-of-$N$
over initial noise \cite{ma2025inferencescaling}, or particle steering
\cite{singhal2025fk}. Only the optimiser changes; the objective and the
evaluation are the same.

\subsection{A reader that needs no download}
\label{sec:reader}

Our in-loop reader is intended to be the dual-channel model of
\citet{sun2025multiorder}, the one model reported by
\citet{rosario2026illusory} to show the human direction under simulated
saccades. Because its weights are gated and its gradient behaviour cannot be
checked in advance, we also supply a small reader with no learned parameters,
so the method is runnable and testable on its own. It has three stages.

\paragraph{Latency.} Elements at different luminances are registered at
slightly different times \cite{conway2005neural,backus2005illusory}, and a
temporal filter of this kind is the mechanism proposed by
\citet{fermuller2010causal}. We apply a one-tap temporal blend whose weight
depends on local luminance $L$,
\begin{equation}
I^{\text{eff}}_t=(1-a)\,I_t+a\,I_{t-1},
\qquad a=\gamma_{\text{lat}}\,(1-L),
\label{eq:latency}
\end{equation}
so darker regions lag. This stage is necessary rather than decorative: with
$\gamma_{\text{lat}}=0$ the reader measures $\rho=+0.012$ on a
peripheral-drift pattern at $192\times192$, against $+0.100$ with
$\gamma_{\text{lat}}=0.75$. Without it there is no illusion to detect.

\paragraph{Quadrature filtering and phase velocity.} For orientation
$\bn_\phi=(\cos\phi,\sin\phi)$ and spatial frequency $f$, we build a Gabor pair
with a Gaussian envelope $g$,
\begin{equation}
G_e(\bu)=g(\bu)\cos\!\left(2\pi f\,\bu\!\cdot\!\bn_\phi\right),
\qquad
G_o(\bu)=g(\bu)\sin\!\left(2\pi f\,\bu\!\cdot\!\bn_\phi\right),
\label{eq:gabor}
\end{equation}
each with its mean subtracted so a flat region gives no response. Writing
$e_t,o_t$ for the responses and $\varphi_t=\operatorname{atan2}(o_t,e_t)$, the
component velocity along $\bn_\phi$ follows the phase-based construction of
\citet{fleet1990phase}. For a grating $I(\bp)=\cos(2\pi f\,\bp\!\cdot\!\bn_\phi)$
shifted by $\bd$, equation \eqref{eq:warp} gives
$I(\bp-\bd)=\cos(2\pi f\,\bp\!\cdot\!\bn_\phi-2\pi f\,\bd\!\cdot\!\bn_\phi)$,
so the true phase changes by $-2\pi f(\bd\!\cdot\!\bn_\phi)$. The convolution in
standard deep-learning libraries is cross-correlation, which negates the
measured phase, and the recovered velocity is therefore
\begin{equation}
v_\phi=\frac{\Delta\varphi}{2\pi f},
\qquad
\Delta\varphi=\operatorname{atan2}\bigl(\sin(\varphi_1-\varphi_0),\,
                                        \cos(\varphi_1-\varphi_0)\bigr),
\label{eq:phase}
\end{equation}
with the $\operatorname{atan2}$ form wrapping the difference to $(-\pi,\pi]$
differentiably. The sign in \eqref{eq:phase} is easy to get wrong, and getting
it wrong reverses every result while leaving all of them plausible; we test it
directly by shifting a texture by a known amount.

\paragraph{Combination.} Each orientation constrains only the velocity
component along $\bn_\phi$, the aperture problem. We combine the constraints by
weighted least squares, with weights $w_\phi=A_0A_1$ given by the product of
the quadrature amplitudes, and Tikhonov damping $\lambda$ so that untextured
regions do not produce large velocities:
\begin{equation}
\bA=\sum_{\phi,f}w_\phi\,\bn_\phi\bn_\phi^{\!\top},
\quad
\bb=\sum_{\phi,f}w_\phi\,v_\phi\,\bn_\phi,
\quad
\bv=(\bA+\lambda\mathbf I)^{-1}\bb .
\label{eq:lstsq}
\end{equation}
This is the standard intersection-of-constraints solution; the recurrent
attention stage of \citet{sun2025multiorder} plays the corresponding role in
the model we use when it is available.

\paragraph{An honest gap.} Equation \eqref{eq:latency} makes lag decrease with
luminance. \citet{conway2005neural} instead report that response latency
depends on \emph{contrast relative to the mean}: black and white elements peak
10--20\,ms before the intermediate greys. We implemented that version as well.
It produces a stronger separation between a drift pattern and its
luminance-reversed twin (0.28 against 0.20 in $\rho$), but with the opposite
sign. We do not resolve which is right here; both are simplifications of a
mechanism that is still debated
\cite{conway2005neural,backus2005illusory,fermuller2010causal}. The practical
consequence is that a reader's direction convention must be calibrated once
against a stimulus whose human direction is known, and that no claim in this
paper rests on a single reader. It is also the clearest argument we can give
for the protocol in Section~\ref{sec:controls}.

\subsection{Controls}
\label{sec:controls}

Three checks are built into the evaluation.

\emph{Held-out readers.} The reader used in \eqref{eq:loss} never appears in
reported numbers. Reported scores come from readers of a different family
\cite{wang2024searaft,lotter2017prednet}, which are never differentiated
through. This is enforced in code rather than by convention.

\emph{Signed control.} Flipping $\omega$ in \eqref{eq:rotation} must flip the
sign of the measured $\rho$. A method that produces motion energy without
controlling its direction passes the unsigned test and fails this one.

\emph{Negative control.} A transform that removes the percept for a human
viewer must drive $\rho$ toward zero. Our first attempt was a single global
permutation of the quantised luminance levels, which preserves the layout and
the set of luminances present. It does not work, and the way it fails is
informative: a global permutation still leaves a coherent asymmetric sequence,
just a different one, so it yields a different illusion rather than none. The
measured $\rho$ rose from $+0.100$ to $+0.390$. What works is permuting the
levels independently within local blocks, which keeps the local histogram and
the contrast but destroys the ordering across the pattern. We retain the global
permutation as a \emph{positive} control, where a high score is the expected
outcome.

\subsection{Stimuli and reproducibility}

Calibration and initialisation use procedurally generated drift patterns rather
than published illusions, which are copyrighted. Elements are laid out in
concentric rings with the asymmetric luminance sequence of
\citet{kitaoka2003snakes}, at an angular repeat count that scales with
resolution. The last point matters: with a fixed repeat count the pattern
aliases at small sizes and the illusion disappears without warning. Measured
$\rho$ on the same pattern rendered at three sizes was $+0.003$ at 96 pixels,
$+0.051$ at 128, and $+0.100$ at 192, with the jitter step scaled with
resolution. Since a stage-I pixel-space diffusion
model generates at $64\times64$, this determines where guidance can be applied,
and we treat the guidance stage as a setting to be checked rather than assumed.

\bibliographystyle{plainnat}
\bibliography{refs}

\end{document}
